\section{Indicator Oracle}

This section exports the compiled proof of the indicator oracle
\(\Uindic\).  The Lean declarations are
\Lean{indicatorOracleImage}, \Lean{indicatorOracleMatrix},
\Lean{indicatorOracleImage\_systemVal\_preserved},
\Lean{indicatorOracleImage\_isBulk\_preserved},
\Lean{indicatorOracleImage\_self\_inverse},
\Lean{indicatorOracleImage\_bijective}, and
\Lean{indicatorOracleMatrix\_is\_permutation}.

\subsection{Definitions}

Let \(n=p.n\) and
\[
  \operatorname{indPos}=1+2n.
\]
For a basis index \(j\), define
\[
  \operatorname{sys}(j)=
  (j\gg 1)\mathbin{\&}(2^n-1).
\]
Let
\[
  \operatorname{bulk}(j)=
  \begin{cases}
  1, & 2\le \operatorname{sys}(j)\le 2^n-3,\\
  0, & \text{otherwise}.
  \end{cases}
\]
The image function is
\[
  I(j)=j\oplus(\operatorname{bulk}(j)\ll \operatorname{indPos}).
\]
The matrix is
\[
  M_{i,j}=
  \begin{cases}
  1, & i=I(j),\\
  0, & i\ne I(j).
  \end{cases}
\]

\subsection{Proof}

The bit \(\operatorname{indPos}\) is above the system register, so toggling it
does not change the extracted system value:
\[
  \operatorname{sys}(I(j))=\operatorname{sys}(j).
\]
This is \Lean{indicatorOracleImage\_systemVal\_preserved}.  Since the bulk
predicate depends only on the system value,
\[
  \operatorname{bulk}(I(j))=\operatorname{bulk}(j),
\]
which is \Lean{indicatorOracleImage\_isBulk\_preserved}.

Therefore
\[
  I(I(j))
  =j\oplus b\ll\operatorname{indPos}\oplus b\ll\operatorname{indPos}
  =j,
\]
where \(b=\operatorname{bulk}(j)\).  This is
\Lean{indicatorOracleImage\_self\_inverse}.

A self-inverse map is injective.  On the finite basis domain it is surjective
because \(I(I(y))=y\).  Lean packages these two facts as
\Lean{indicatorOracleImage\_bijective}.  The entry characterization
\Lean{indicatorOracleMatrix\_eq\_image} then gives exactly one \(1\)-entry in
each row and each column.  This is
\Lean{indicatorOracleMatrix\_is\_permutation}.

Consequently \Lean{oneTermRobinGate\_U\_indic} has
\Lean{unitary.proved = true}.
